% GENERATED FILE. Edit canonical lessons, then run npm run generate:books.
% canonical-source-hash: fnv1a64:dcda597d1edb88b9
% canonical-lessons: MW-C209-baniyo, MW-C209-gahak, MW-C209-seth, MW-C209-saudo, MW-C209-mangavno, MW-R209-review-204-209

\chapter{Shopkeeper and Customer}
\label{ch:mw-shopkeeper-and-customer}
\hlchaptermodality{209}

\begin{chapteropening}
\textbf{By the end of this chapter:} \emph{I can say who sells, who buys and what I order.}

Complete the last lesson of chapter 209: I can say who sells, who buys and what I order.
\end{chapteropening}

\section[\texorpdfstring{\mw{बाणियो} (bāṇiyo)}{बाणियो (bāṇiyo)}]{\mw{बाणियो} (bāṇiyo) — a trader; a shopkeeper}

\label{lesson:MW-C209-baniyo}



\begin{quote}
Before the new one: say the Marwadi for a court of law, then the Marwadi for to come open.
\end{quote}

\subsection*{You'll want to know: \mw{बाणियो}}
\textbf{\mw{बाणियो}} — \emph{bāṇiyo} — \textquotedblleft{}a trader; a shopkeeper\textquotedblright{}.

\begin{quote}
\mw{बाणियो} \mw{दिनूगै} \mw{दुकान} \mw{खोलै} \mw{है।}
\end{quote}

\begin{quote}
\emph{bāṇiyo dinūgai dukān kholai hai.}
\end{quote}

\begin{quote}
The trader opens the shop early in the morning.
\end{quote}

A masculine noun.

\begin{cousinweb}
\textbf{\mw{बाणियो}} (\emph{bāṇiyo}) is inherited from Sanskrit \emph{vāṇija}, \textquotedblleft{}trader\textquotedblright{}; CDIAL lists Marwari \emph{bā̃ṇyõ}, \textquotedblleft{}trader\textquotedblright{} (CDIAL 11484).
\end{cousinweb}

\begin{grammarlens}[title={What you've built}]
One more word for this chapter.
\end{grammarlens}

\subsection*{Your turn}
\begin{itemize}
\raggedright
  \item \emph{Say it:} \emph{bāṇiyo}
  \item \emph{Say it:} \emph{bāṇiyo}, once more
  \item \emph{Say it:} \emph{bāṇiyo dinūgai dukān kholai hai.} — The trader opens the shop early in the morning.
  \item \emph{Recall:} say the Marwadi for a court of law, then the Marwadi for to come open, then say \emph{bāṇiyo} again
\end{itemize}

\begin{tcolorbox}[breakable,colback=teal!4,colframe=teal!35!black,title={Before you move on}]
What does \mw{बाणियो} mean? (\textquotedblleft{}A trader; a shopkeeper\textquotedblright{}.) Say it once more.
\end{tcolorbox}
\section[\texorpdfstring{\mw{गाहक} (gāhak)}{गाहक (gāhak)}]{\mw{गाहक} (gāhak) — a customer}

\label{lesson:MW-C209-gahak}



\begin{quote}
Before the new one: say the Marwadi for to come open, then the Marwadi for a trader.
\end{quote}

\subsection*{You'll want to know: \mw{गाहक}}
\textbf{\mw{गाहक}} — \emph{gāhak} — \textquotedblleft{}a customer\textquotedblright{}.

\begin{quote}
\mw{गाहक} \mw{भाव} \mw{पूछै} \mw{है।}
\end{quote}

\begin{quote}
\emph{gāhak bhāv pūchhai hai.}
\end{quote}

\begin{quote}
The customer asks the price.
\end{quote}

A masculine or feminine noun, by the person meant.

\begin{cousinweb}
\textbf{\mw{गाहक}} (\emph{gāhak}) is a tadbhava of Sanskrit \emph{grāhaka}, \textquotedblleft{}one who takes; a purchaser\textquotedblright{} (Monier-Williams), with the \emph{r} lost.
\end{cousinweb}

\begin{grammarlens}[title={What you've built}]
One more word for this chapter.
\end{grammarlens}

\subsection*{Your turn}
\begin{itemize}
\raggedright
  \item \emph{Say it:} \emph{gāhak}
  \item \emph{Say it:} \emph{gāhak}, once more
  \item \emph{Say it:} \emph{gāhak bhāv pūchhai hai.} — The customer asks the price.
  \item \emph{Recall:} say the Marwadi for to come open, then the Marwadi for a trader, then say \emph{gāhak} again
\end{itemize}

\begin{tcolorbox}[breakable,colback=teal!4,colframe=teal!35!black,title={Before you move on}]
What does \mw{गाहक} mean? (\textquotedblleft{}A customer\textquotedblright{}.) Say it once more.
\end{tcolorbox}
\section[\texorpdfstring{\mw{सेठ} (seṭh)}{सेठ (seṭh)}]{\mw{सेठ} (seṭh) — a merchant}

\label{lesson:MW-C209-seth}



\begin{quote}
Before the new one: say the Marwadi for a trader, then the Marwadi for a customer.
\end{quote}

\subsection*{You'll want to know: \mw{सेठ}}
\textbf{\mw{सेठ}} — \emph{seṭh} — \textquotedblleft{}a merchant\textquotedblright{}.

\begin{quote}
\mw{सेठ} \mw{हवेली} \mw{मांय} \mw{रैवै} \mw{है।}
\end{quote}

\begin{quote}
\emph{seṭh havelī mā̃y raivai hai.}
\end{quote}

\begin{quote}
The merchant lives in a mansion.
\end{quote}

A masculine noun.

\begin{cousinweb}
\textbf{\mw{सेठ}} (\emph{seṭh}) is inherited from Sanskrit \emph{śreṣṭhin}, \textquotedblleft{}a distinguished man; the head of a guild\textquotedblright{}. CDIAL gives no Marwari form in this entry, but lists Hindi \emph{seṭh}, \textquotedblleft{}merchant, banker\textquotedblright{} (CDIAL 12726).
\end{cousinweb}

\begin{grammarlens}[title={What you've built}]
One more word for this chapter.
\end{grammarlens}

\subsection*{Your turn}
\begin{itemize}
\raggedright
  \item \emph{Say it:} \emph{seṭh}
  \item \emph{Say it:} \emph{seṭh}, once more
  \item \emph{Say it:} \emph{seṭh havelī mā̃y raivai hai.} — The merchant lives in a mansion.
  \item \emph{Recall:} say the Marwadi for a trader, then the Marwadi for a customer, then say \emph{seṭh} again
\end{itemize}

\begin{tcolorbox}[breakable,colback=teal!4,colframe=teal!35!black,title={Before you move on}]
What does \mw{सेठ} mean? (\textquotedblleft{}A merchant\textquotedblright{}.) Say it once more.
\end{tcolorbox}
\section[\texorpdfstring{\mw{सौदो} (saudo)}{सौदो (saudo)}]{\mw{सौदो} (saudo) — a deal; shopping}

\label{lesson:MW-C209-saudo}



\begin{quote}
Before the new one: say the Marwadi for a customer, then the Marwadi for a merchant.
\end{quote}

\subsection*{You'll want to know: \mw{सौदो}}
\textbf{\mw{सौदो}} — \emph{saudo} — \textquotedblleft{}a deal; shopping\textquotedblright{}.

\begin{quote}
\mw{म्हैं} \mw{बाजार} \mw{मांय} \mw{सौदो} \mw{कियो।}
\end{quote}

\begin{quote}
\emph{mhaĩ bājār mā̃y saudo kiyo.}
\end{quote}

\begin{quote}
I did the shopping at the market.
\end{quote}

A masculine noun.

\begin{cousinweb}
\textbf{\mw{सौदो}} (\emph{saudo}) is usually traced to Persian \emph{saudā}, \textquotedblleft{}trade, a bargain\textquotedblright{}.
\end{cousinweb}

\begin{grammarlens}[title={What you've built}]
One more word for this chapter.
\end{grammarlens}

\subsection*{Your turn}
\begin{itemize}
\raggedright
  \item \emph{Say it:} \emph{saudo}
  \item \emph{Say it:} \emph{saudo}, once more
  \item \emph{Say it:} \emph{mhaĩ bājār mā̃y saudo kiyo.} — I did the shopping at the market.
  \item \emph{Recall:} say the Marwadi for a customer, then the Marwadi for a merchant, then say \emph{saudo} again
\end{itemize}

\begin{tcolorbox}[breakable,colback=teal!4,colframe=teal!35!black,title={Before you move on}]
What does \mw{सौदो} mean? (\textquotedblleft{}A deal; shopping\textquotedblright{}.) Say it once more.
\end{tcolorbox}
\section[\texorpdfstring{\mw{मंगावणो} (maṅgāvṇo)}{मंगावणो (maṅgāvṇo)}]{\mw{मंगावणो} (maṅgāvṇo) — to order; to have brought}

\label{lesson:MW-C209-mangavno}



\begin{quote}
Before the new one: say the Marwadi for a merchant, then the Marwadi for a deal.
\end{quote}

\subsection*{You'll want to know: \mw{मंगावणो}}
\textbf{\mw{मंगावणो}} — \emph{maṅgāvṇo} — \textquotedblleft{}to order; to have brought\textquotedblright{}.

\begin{quote}
\mw{हूं} \mw{होटल} \mw{सूं} \mw{रोटी} \mw{मंगाऊं} \mw{हूं।}
\end{quote}

\begin{quote}
\emph{hū̃ hoṭal sū̃ roṭī maṅgāū̃ hū̃.}
\end{quote}

\begin{quote}
I order flatbread from the hotel.
\end{quote}

Past (I …): \mw{म्हैं} \mw{चाय} \mw{मंगाई}, \emph{mhaĩ chāy maṅgāī}.

\begin{cousinweb}
\textbf{\mw{मंगावणो}} (\emph{maṅgāvṇo}) is the causative of \textbf{\mw{मांगणो}}, \textquotedblleft{}to ask for\textquotedblright{}: to ask for a thing to be brought. Both are inherited from Sanskrit \emph{mārgati}, \textquotedblleft{}seeks, requests\textquotedblright{}; CDIAL lists Old Marwari \emph{māgaï}, \textquotedblleft{}asks for\textquotedblright{} (CDIAL 10074).
\end{cousinweb}

\begin{grammarlens}[title={What you've built}]
One more word for this chapter.
\end{grammarlens}

\subsection*{Your turn}
\begin{itemize}
\raggedright
  \item \emph{Say it:} \emph{maṅgāvṇo}
  \item \emph{Say it:} \emph{maṅgāvṇo}, once more
  \item \emph{Say it:} \emph{hū̃ hoṭal sū̃ roṭī maṅgāū̃ hū̃.} — I order flatbread from the hotel.
  \item \emph{Recall:} say the Marwadi for a merchant, then the Marwadi for a deal, then say \emph{maṅgāvṇo} again
\end{itemize}

\begin{tcolorbox}[breakable,colback=teal!4,colframe=teal!35!black,title={Before you move on}]
What does \mw{मंगावणो} mean? (\textquotedblleft{}To order; to have brought\textquotedblright{}.) Say it once more.
\end{tcolorbox}
\section[(dialogue)]{Looking back: from Out and About to Shopkeeper and Customer}

\label{lesson:MW-R209-review-204-209}



\begin{quote}
Say the Marwadi for to stroll and a crowd.
\end{quote}

\subsection*{Your turn}
Say these aloud:

\begin{itemize}
\raggedright
  \item to stroll, a crowd, an open ground, an edge, a map
  \item a ticket, to miss, a fare, to get stuck, a bus conductor
  \item third, fourth, upstairs, on the other side, separate
  \item a garden, a water tank, a checkpoint, a big gate, a tower
  \item a hospital, a sub-district office, a factory, a court of law, to come open
  \item a trader, a customer, a merchant, a deal, to order
\end{itemize}

\begin{tcolorbox}[breakable,colback=teal!4,colframe=teal!35!black,title={Before you move on}]
To stroll? (\textbf{\mw{घूमणो}}, \emph{ghūmṇo}.) A crowd? (\textbf{\mw{भीड़}}, \emph{bhīṛ}.) An open ground? (\textbf{\mw{मैदान}}, \emph{maidān}.) An edge? (\textbf{\mw{किनारो}}, \emph{kināro}.) A map? (\textbf{\mw{नक्सो}}, \emph{nakso}.) A ticket? (\textbf{\mw{टिकट}}, \emph{ṭikaṭ}.) To miss? (\textbf{\mw{चूकणो}}, \emph{chūkṇo}.) A fare? (\textbf{\mw{भाड़ो}}, \emph{bhāṛo}.) To get stuck? (\textbf{\mw{अटकणो}}, \emph{aṭakṇo}.) A bus conductor? (\textbf{\mw{कंडक्टर}}, \emph{kaṇḍakṭar}.) Third? (\textbf{\mw{तीजो}}, \emph{tījo}.) Fourth? (\textbf{\mw{चौथो}}, \emph{chautho}.) Upstairs? (\textbf{\mw{ऊपर}}, \emph{ūpar}.) On the other side? (\textbf{\mw{पार}}, \emph{pār}.) Separate? (\textbf{\mw{अळगो}}, \emph{aḷgo}.) A garden? (\textbf{\mw{बगीचो}}, \emph{bagīcho}.) A water tank? (\textbf{\mw{कुंड}}, \emph{kuṇḍ}.) A checkpoint? (\textbf{\mw{नाको}}, \emph{nāko}.) A big gate? (\textbf{\mw{फाटक}}, \emph{phāṭak}.) A tower? (\textbf{\mw{बुरज}}, \emph{buraj}.) A hospital? (\textbf{\mw{अस्पताल}}, \emph{aspatāl}.) A sub-district office? (\textbf{\mw{तहसील}}, \emph{tahsīl}.) A factory? (\textbf{\mw{कारखानो}}, \emph{kārkhāno}.) A court of law? (\textbf{\mw{अदालत}}, \emph{adālat}.) To come open? (\textbf{\mw{खुलणो}}, \emph{khulṇo}.) A trader? (\textbf{\mw{बाणियो}}, \emph{bāṇiyo}.) A customer? (\textbf{\mw{गाहक}}, \emph{gāhak}.) A merchant? (\textbf{\mw{सेठ}}, \emph{seṭh}.) A deal? (\textbf{\mw{सौदो}}, \emph{saudo}.) To order? (\textbf{\mw{मंगावणो}}, \emph{maṅgāvṇo}.)
\end{tcolorbox}
